\section{The Minimum Viable Deployment}
\label{sec:vision}

This section states what at a minimum has to be added to a vehicle for deploying
the self-regulating cars vision.

\noindent\textbf{Two sensing layers, and two observation spaces.}
The deployment carries two decisions, and each reads its own sensor. The
advisory layer decides a desired speed for the super-segment the vehicle is on.
The safety layer decides whether that speed may be shown to the driver.
Table~\ref{tab:layers} states the split, and the two paragraphs after it state
what each layer observes.

\begin{table}[t]
\centering
\caption{The two layers, the sensor each reads, and the decision each makes.
Neither sensor can supply the other layer's inputs.}
\label{tab:layers}
\small
\resizebox{\columnwidth}{!}{%
\begin{tabular}{p{0.15\linewidth}p{0.19\linewidth}p{0.26\linewidth}p{0.19\linewidth}}
\toprule
\textbf{Layer} & \textbf{Sensor} & \textbf{Produces} & \textbf{Decides} \\
\midrule
Advisory & Traffic API, per super-segment & Aggregate speed, free flow, jam factor & The desired speed \\
Safety & Camera, GPS, IMU & Leader distance, closing speed, ego state & Whether to allow it \\
\bottomrule
\end{tabular}
}
\end{table}

\noindent\textbf{The advisory observation space.}
The advisory layer reads the whole network at once, as a fixed partition into
super-segments, with five fields for each. 
Three fields come from the traffic feed, which in this deployment is the HERE
Traffic API~\cite{here_traffic_api}, and two from the map. \texttt{speed}
is the segment's current aggregate speed, which is the quantity being
regulated. \texttt{free\_flow} is the speed the segment carries when it is
uncongested, so the two together state how far below its own capacity the
segment is running. \texttt{jam\_factor} is a congestion index on a 0 to 10
scale, which summarises the same shortfall on a bounded scale the feed
publishes directly. \texttt{lanes} turns a per-segment reading into a per-lane
one, and \texttt{length\_km} says how much road that reading describes. The
feed publishes \texttt{confidence} and \texttt{traversability} as well, but
neither is read.

The action is one of three speed fractions per super-segment. The base it
multiplies is that segment's static speed limit from the map, and it is
deliberately not the feed's free-flow reading, because a live reading used for
the base as well as for the observation would enter the loop twice. The driver
sees one number, which is the row for the super-segment the vehicle's own
position currently falls on.

\noindent\textbf{The safety observation space.}
The safety layer reads the road immediately around the vehicle. The forward
camera supplies the distance to the vehicle ahead and the rate at which that
distance is closing. GPS supplies ego speed, and the slope of that speed over
a short window supplies ego acceleration. The tracked vehicles in range supply
a local density and a local mean speed. Those are the quantities needed to
decide that a recommended speed is unsafe or obstructive at this moment, which
is a judgement about the road inside the camera's range rather than about the
network.

\noindent\textbf{The reasoning for the split.}
A traffic API reports aggregate quantities per road segment and contains no
individual vehicles, so it cannot report a leader distance, a follower gap or
a lane distribution. A forward camera is the opposite. It measures the vehicle
ahead and reports nothing about the road beyond its own horizon. Neither
sensor can do the other's job, hence a clear split design.

\noindent\textbf{Why not aggregate from local sensing.}
Another design could be local sensing, where a vehicle reconstructs segment-level state from its own
sensors and those of its neighbours. That needs either vehicle-to-vehicle
exchange, or a server that aggregates. The
second option is exactly the deployment cost we want to avoid. Querying a traffic API
per vehicle is simpler to operate than having vehicles talk to each other. Also, for 
local sensing to create a rich segment state, a sufficient number of compliant vehicles is 
required. Such a system can not be tested with a unit deployment. In the design reported 
here the per-vehicle quantities stay where they are measurable, and they feed the safety layer only.

\noindent\textbf{Decentralized execution.}
The policy's input is the whole network's state, and the policy is a
deterministic argmax over three speed fractions. Every vehicle that holds the
same weights, reads the same map and queries the same feed computes the same
action for the same super-segment. Decentralized execution therefore reproduces
centralized execution by construction. What does differ on a road
against the simulation is staleness. A response carries the time it was
received, which this system measures. It does not carry the time at which the
conditions it describes held, and nothing in the response reports that. The query is small. O(10) super-segments at five
features is O(100) numbers per decision, and the decision interval is 60\,s. The
feed is queried by corridor or by circle, so one request covers the whole extent
and the number of super-segments does not change the number of requests.

\noindent\textbf{What the vehicle must therefore carry.}
Five things. (1) A network definition that names the super-segments and
their edges. (2) A traffic-feed query bounded by that network's extent. (3) A policy
trained on that network. (4) A local sensor set for the safety layer. (5) A display. 
There is no connection to the
throttle, the brakes, the steering or the vehicle bus, and no server inside the
control loop.
